\chapter{Finiteness of cohomology of line bundles on a proper curve (S2)}
\label{chap:Cohomology_SerreFiniteness}

% Chapter-local shorthands (report to plan agent for promotion to
% macros/common.tex if shared across chapters).
\providecommand{\cechHI}{\mathtt{cechH1}}
\providecommand{\FiniteDimensional}{\mathrm{FiniteDimensional}}

% SOURCE: Hartshorne, \emph{Algebraic Geometry}, GTM 52 (1977): II.5,
% Theorem~5.19 (finite generation of global sections of a coherent sheaf
% on a projective scheme); III.2, Theorem~2.7 (Grothendieck vanishing
% above the dimension); III.5, Theorem~5.2 + Remark~5.2.1 (Serre's
% coherent-cohomology finiteness on a projective scheme). All verbatim
% quotes in the blocks below were read off the rendered scanned page
% images of references/hartshorne-algebraic-geometry.pdf (PDF page 139 =
% doc p.122 for II.5.19; PDF page 225 = doc p.208 for III.2.7; PDF page
% 228 = doc p.228 for III.5.2) and transcribed character-by-character;
% the PDF has no text layer.

\section{Setup and motivation}

This chapter supplies the \texttt{S2} finiteness obligation of the
genus-\(0\) Riemann--Roch arm: \emph{for every Weil divisor \(D\) on a
smooth proper geometrically irreducible curve \(C / \bar k\), the
sheaf-cohomology groups \(H^i(C, \mathcal O_C(D))\) are
finite-dimensional \(\bar k\)-vector spaces}. These finiteness facts are
exactly what makes the \(\ell\)-invariant \(\ell(D) = \dim_{\bar k}
H^0(C, \mathcal O_C(D))\) (\cref{def:l_invariant}) and the Euler
characteristic \(\chi(\mathcal O_C(D)) = \dim_{\bar k} H^0 -
\dim_{\bar k} H^1\) (\cref{def:eulerChar_curve}) honest integers, and
they discharge the four \texttt{Module.Finite} side-conditions consumed
by the \(\chi\)-inductive step
\cref{lem:euler_char_sheafOf_succ} of \cref{chap:RiemannRoch_RRFormula}.

\paragraph{Strategy of the chapter (the SES-induction route).} Rather
than formalising Serre's general coherent-cohomology finiteness theorem
(the projective-dévissage statement), the chapter reduces \emph{all} the
finiteness it needs to a single base \emph{computation} and one
structural induction, using supporting results proved in related
chapters:
\begin{enumerate}
  \item the closed-point skyscraper \(k(P)\) has finite cohomology,
    read off from the skyscraper \(H^0\)/\(H^1\) computations
    (\cref{sec:SerreFiniteness_skyscraper});
  \item finiteness is a \emph{2-of-3} property along a short exact
    sequence of sheaves, via the long exact cohomology sequence
    (\cref{lem:HModule_finite_of_ses});
  \item the base case is the structure sheaf \(\mathcal O_C\)
    (\cref{lem:HModule_structureSheaf_finite}); its \(H^1\) finiteness
    is computed \emph{directly} from a two-affine Čech cover
    (\cref{lem:HModule_structureSheaf_H1_finite}): writing
    \(C = U \cup V\) with \(U, V\) affine opens, one has
    \(H^1(C, \mathcal O_C) \cong
    \operatorname{coker}\bigl(\mathcal O_C(U) \oplus \mathcal O_C(V)
    \to \mathcal O_C(U \cap V)\bigr)\), and this explicit cokernel is a
    finite-dimensional \(\bar k\)-vector space;
  \item every line bundle \(\mathcal O_C(D)\) is reached from
    \(\mathcal O_C\) by repeatedly adding and removing single closed
    points, along the closed-point short exact sequence
    \cref{lem:sheafOf_ses_single_add}; finiteness propagates by the
    2-of-3 property (\cref{thm:HModule_lineBundle_finite}).
\end{enumerate}
This keeps the entire \texttt{S2} obligation reduced to one base
computation --- the \(H^1\) of \(\mathcal O_C\) read off a two-affine
Čech cover --- together with one structural induction, and avoids the
proper-pushforward / projective dévissage machinery not available at the
\(\HModule \bar k\) cohomology level.

\paragraph{Standing hypotheses.} Throughout, \(\bar k\) is an
algebraically closed field and \(C\) is a smooth proper geometrically
irreducible curve over \(\bar k\).
Cohomology is the \(\Module \bar k\)-flavoured \(\Ext\)
cohomology \(\HModule\;\bar k\;\mathcal F\;i\)
(\cref{def:Scheme_HModule}); we write
\(H^i(C, \mathcal F) := \HModule\;\bar k\;\mathcal F\;i\). For a sheaf
\(\mathcal F\) of \(\bar k\)-modules, ``\(\mathcal F\) has finite
cohomology'' abbreviates: \(H^i(C, \mathcal F)\) is a finite
\(\bar k\)-module (\(\Module.\mathtt{Finite}\;\bar k\)) for every
\(i\); since \(\bar k\) is a field this is the same as each
\(H^i(C, \mathcal F)\) being a finite-dimensional \(\bar k\)-vector
space.


\section{The degree bound}
\label{sec:SerreFiniteness_inputs}

The chapter leans on one external statement: Grothendieck's vanishing
theorem, used to bound the relevant degrees to \(i \in \{0, 1\}\). The
single base obligation it leaves --- finiteness of \(H^1(\mathcal O_C)\)
--- is then discharged \emph{internally} by the two-affine Čech
computation of \cref{sec:SerreFiniteness_cech}, not admitted as a gap.

\begin{lemma}
[Grothendieck vanishing bounds the relevant degrees]
  \label{lem:grothendieck_vanishing_curve}
  % SOURCE: [Hartshorne], III.2, Theorem~2.7 (Grothendieck) (read from
  % references/hartshorne-algebraic-geometry.pdf, PDF page 225 = doc
  % p.208).
  % SOURCE QUOTE: "Theorem 2.7 (Grothendieck [1]). \emph{Let $X$ be a
  % noetherian topological space of dimension $n$. Then for all $i > n$
  % and all sheaves of abelian groups $\mathscr{F}$ on $X$, we have
  % $H^i(X, \mathscr{F}) = 0$.}"
  \textit{Source: Hartshorne, III.2, Theorem~2.7, p.~208.}
  Let \(C\) be a smooth proper geometrically irreducible curve over
  \(\bar k\); its underlying topological space is noetherian of
  dimension \(1\). Then, for every sheaf \(\mathcal F\) of
  \(\bar k\)-modules on \(C\) and every \(i \ge 2\),
  \[
    H^i(C, \mathcal F) \;=\; 0.
  \]
  In particular \(H^i(C, \mathcal F)\) is a finite (indeed zero)
  \(\bar k\)-module for all \(i \ge 2\), so to prove that \(\mathcal F\)
  has finite cohomology it suffices to treat the two degrees
  \(i \in \{0, 1\}\).
\end{lemma}

\begin{proof}
  \uses{def:Scheme_HModule}
  The space underlying \(C\) is noetherian of Krull dimension \(1\)
  (an integral curve), so Grothendieck's vanishing theorem
  (Hartshorne~III.2.7) with \(n = 1\) gives
  \(H^i(C, \mathcal F) = 0\) for all \(i > 1\), i.e.\ all \(i \ge 2\).
  The zero module is a finite \(\bar k\)-module, so the only degrees
  carrying content for the finiteness question are \(i = 0\) and
  \(i = 1\).
\end{proof}

\section{Finiteness for the closed-point skyscraper}
\label{sec:SerreFiniteness_skyscraper}

\begin{lemma}
[Finiteness of the cohomology of a closed-point skyscraper]
  \label{lem:HModule_skyscraper_finite}
  \lean{AlgebraicGeometry.Scheme.HModule_skyscraper_finite}
  \uses{lem:H0_skyscraperSheaf_finrank_eq_one,
        lem:H1_skyscraperSheaf_finrank_eq_zero_main,
        lem:grothendieck_vanishing_curve, def:Scheme_HModule}
  Let \(C\) be a smooth proper geometrically irreducible curve over
  \(\bar k\) and let \(P \in C\) be a closed point. Write
  \[
    k(P) \;:=\; \mathtt{skyscraperSheaf}\;P\;
      (\Module.\mathtt{of}\;\bar k\;\bar k)
  \]
  for the closed-point skyscraper sheaf. Then \(k(P)\) has finite
  cohomology: \(H^i(C, k(P))\) is a finite-dimensional \(\bar k\)-vector
  space for every \(i\). Concretely \(H^0(C, k(P)) \cong \bar k\) (so
  \(\dim_{\bar k} H^0 = 1\)), \(H^1(C, k(P)) = 0\), and
  \(H^i(C, k(P)) = 0\) for \(i \ge 2\).

\end{lemma}

\begin{proof}
  \uses{lem:H0_skyscraperSheaf_finrank_eq_one,
        lem:H1_skyscraperSheaf_finrank_eq_zero_main,
        lem:grothendieck_vanishing_curve}
  We check the three degree-regimes.

  \emph{Degree \(0\).} The skyscraper evaluation
  \cref{lem:H0_skyscraperSheaf_finrank_eq_one} of
  \cref{chap:RiemannRoch_RRFormula} gives
  \(\dim_{\bar k} H^0(C, k(P)) = 1\): a global section of the skyscraper
  at \(P\) is precisely an element of its stalk \(\bar k\), since
  \(P\) lies in every open whose section group is non-trivial. Hence
  \(H^0(C, k(P)) \cong \bar k\) is one-dimensional, in particular a
  finite \(\bar k\)-module.

  \emph{Degree \(1\).} The headline skyscraper vanishing
  \cref{lem:H1_skyscraperSheaf_finrank_eq_zero_main} of
  \cref{chap:RR_H1Vanishing} gives \(\dim_{\bar k} H^1(C, k(P)) = 0\),
  i.e.\ \(H^1(C, k(P)) = 0\); the zero module is finite.

  \emph{Degrees \(i \ge 2\).} By
  \cref{lem:grothendieck_vanishing_curve} the cohomology of any sheaf of
  \(\bar k\)-modules on the one-dimensional space \(C\) vanishes in
  degree \(\ge 2\), so \(H^i(C, k(P)) = 0\) there as well.

  In every degree \(H^i(C, k(P))\) is therefore a finite-dimensional
  \(\bar k\)-vector space.
\end{proof}

\section{Finiteness is a 2-of-3 property along a short exact sequence}

\begin{lemma}
[Two-out-of-three for finite cohomology along a short exact sequence]
  \label{lem:HModule_finite_of_ses}
  \lean{AlgebraicGeometry.Scheme.HModule_finite_of_ses}
  \uses{def:Scheme_HModule, lem:grothendieck_vanishing_curve}
  Let \(C\) be a smooth proper geometrically irreducible curve over
  \(\bar k\) and let
  \[
    0 \;\longrightarrow\; \mathcal F_1 \;\longrightarrow\;
    \mathcal F_2 \;\longrightarrow\; \mathcal F_3 \;\longrightarrow\; 0
  \]
  be a short exact sequence of sheaves of \(\bar k\)-modules on \(C\).
  If two of the three sheaves have finite cohomology (each \(H^i\) a
  finite-dimensional \(\bar k\)-vector space), then so does the third.

\end{lemma}

\begin{proof}
  \uses{def:Scheme_HModule, lem:grothendieck_vanishing_curve}
  The short exact sequence induces, through the \(\HModule \bar k\)
  cohomology functor (a \(\delta\)-functor on the \(\bar k\)-linear abelian
  sheaf category; the same long exact sequence used for \(\chi\)-additivity
  in \cref{chap:RiemannRoch_RRFormula}), a long exact sequence of
  \(\bar k\)-vector spaces
  \[
    \cdots \to H^{i-1}(C, \mathcal F_3)
    \to H^{i}(C, \mathcal F_1)
    \to H^{i}(C, \mathcal F_2)
    \to H^{i}(C, \mathcal F_3)
    \to H^{i+1}(C, \mathcal F_1) \to \cdots
  \]
  We use the elementary linear-algebra fact: \emph{in an exact sequence
  \(X \to Y \to Z\) of \(\bar k\)-vector spaces with \(X\) and \(Z\)
  finite-dimensional, the middle term \(Y\) is finite-dimensional}
  (the image of \(X\) is finite-dimensional, and \(Y\) modulo that image
  embeds into \(Z\); equivalently the closure of
  \(\Module.\mathtt{Finite}\;\bar k\) under subobjects and extensions in
  the category of \(\bar k\)-modules). Apply it at each degree to the
  appropriate three-term window of the long exact sequence:
  \begin{itemize}
    \item to recover \(\mathcal F_2\) finite from \(\mathcal F_1,
      \mathcal F_3\) finite, read off the window
      \(H^i(\mathcal F_1) \to H^i(\mathcal F_2) \to
      H^i(\mathcal F_3)\), whose two ends are finite-dimensional;
    \item to recover \(\mathcal F_3\) finite from \(\mathcal F_1,
      \mathcal F_2\) finite, read off
      \(H^i(\mathcal F_2) \to H^i(\mathcal F_3) \to
      H^{i+1}(\mathcal F_1)\);
    \item to recover \(\mathcal F_1\) finite from \(\mathcal F_2,
      \mathcal F_3\) finite, read off
      \(H^{i-1}(\mathcal F_3) \to H^i(\mathcal F_1) \to
      H^i(\mathcal F_2)\).
  \end{itemize}
  By \cref{lem:grothendieck_vanishing_curve} only the degrees
  \(i \in \{0, 1\}\) need to be checked (all groups vanish, hence are
  finite, in degree \(\ge 2\)), so the three windows above are applied
  finitely many times. In each case the unknown term is sandwiched
  between two finite-dimensional spaces and is therefore
  finite-dimensional, which is the claimed conclusion.
\end{proof}

\section{The base case via a two-affine Čech cover}
\label{sec:SerreFiniteness_cech}

The single base obligation left by the degree bound is the
finite-dimensionality of \(H^1(C, \mathcal O_C)\). We discharge it
\emph{internally}, by computing that group as the cokernel of an
explicit \(\bar k\)-linear map: a smooth proper curve is covered by two
affine opens, the Čech complex of a two-element cover has only two
terms, and on such an affine cover Čech cohomology agrees with derived
cohomology.

\begin{lemma}
[A smooth proper curve has a two-affine open cover]
  \label{lem:exists_two_affine_cover}
  \lean{AlgebraicGeometry.Scheme.exists_two_affine_open_cover}
  % SOURCE: [Hartshorne], IV.1, Exercise~1.3 (a non-proper curve is
  % affine) together with the IV.1 standing convention that a curve is
  % proper, hence projective (read from
  % references/hartshorne-algebraic-geometry.pdf, PDF page 314 = doc
  % p.297, and PDF page 311 = doc p.294).
  % SOURCE QUOTE: "1.3. Let $X$ be an integral, separated, regular,
  % one-dimensional scheme of finite type over $k$, which is \emph{not}
  % proper over $k$. Then $X$ is affine. [\emph{Hint:} Embed $X$ in a
  % (proper) curve $\bar X$ over $k$, and use (Ex. 1.2) to construct a
  % morphism $f \colon \bar X \to \mathbf{P}^1$ such that
  % $f^{-1}(\mathbf{A}^1) = X$.]"
  % SOURCE QUOTE: "In this chapter we will use the word \emph{curve} to
  % mean a complete, nonsingular curve over an algebraically closed
  % field $k$. In other words (II, \S6), a curve is an integral scheme
  % of dimension 1, proper over $k$, all of whose local rings are
  % regular. ... Such a curve is necessarily projective (II, 6.7)."
  \textit{Source: Hartshorne, IV.1, Exercise~1.3, p.~297, and the IV.1
  curve convention, p.~294.}
  Let \(C\) be a smooth proper geometrically irreducible curve over
  \(\bar k\). Then there exist two distinct closed points
  \(P_0 \ne Q_0 \in C\) such that, setting
  \[
    U \;:=\; C \setminus \{P_0\}, \qquad V \;:=\; C \setminus \{Q_0\},
  \]
  the opens \(U\) and \(V\) are affine (\(\IsAffineOpen\)) and
  \(\{U, V\}\) is an open cover of \(C\), i.e.\ \(U \cup V = C\).
\end{lemma}

\begin{proof}
  Since \(C\) is a positive-dimensional integral curve over the
  algebraically closed field \(\bar k\), it has infinitely many closed
  points; choose two distinct ones \(P_0 \ne Q_0\). Each of
  \(U = C \setminus \{P_0\}\) and \(V = C \setminus \{Q_0\}\) is a
  non-empty open subscheme of \(C\), hence is itself integral,
  separated, regular, one-dimensional, and of finite type over
  \(\bar k\). Neither is proper over \(\bar k\): a proper curve minus a
  closed point is no longer complete. By Hartshorne's Exercise~IV.1.3 a
  scheme with exactly these properties is affine, so \(U\) and \(V\) are
  affine opens. Finally \(\{U, V\}\) covers \(C\): a point other than
  \(P_0\) lies in \(U\), the point \(P_0\) (being distinct from
  \(Q_0\)) lies in \(V\), so every point of \(C\) lies in \(U\) or
  \(V\); hence \(U \cup V = C\).
\end{proof}

\begin{definition}
[First Čech cohomology of the two-affine cover]
  \label{def:cech_H1_two_affine}
  \lean{AlgebraicGeometry.Scheme.cechH1OfTwoAffine}
  \uses{lem:exists_two_affine_cover}
  % SOURCE: [Hartshorne], III.4, the Čech complex
  % $C^\cdot(\mathfrak{U},\mathscr{F})$ and its coboundary, and
  % Example~4.0.3 (a two-open cover gives a two-term complex) (read from
  % references/hartshorne-algebraic-geometry.pdf, PDF page 235 = doc
  % p.218, and PDF page 236 = doc p.219).
  % SOURCE QUOTE: "Now let $\mathscr{F}$ be a sheaf of abelian groups on
  % $X$. We define a complex $C^\cdot(\mathfrak{U},\mathscr{F})$ of
  % abelian groups as follows. For each $p \geq 0$, let
  % $C^p(\mathfrak{U},\mathscr{F}) = \prod_{i_0 < \ldots < i_p}
  % \mathscr{F}(U_{i_0,\ldots,i_p})$. ... We define the co-boundary map
  % $d \colon C^p \to C^{p+1}$ by setting $(d\alpha)_{i_0,\ldots,i_{p+1}}
  % = \sum_{k=0}^{p+1} (-1)^k
  % \alpha_{i_0,\ldots,\hat{\imath_k},\ldots,i_{p+1}}
  % |_{U_{i_0,\ldots,i_{p+1}}}$."
  % SOURCE QUOTE: "Then the Čech complex has only two terms: $C^0 =
  % \Gamma(U,\Omega) \times \Gamma(V,\Omega)$ $C^1 = \Gamma(U \cap
  % V,\Omega)$".
  \textit{Source: Hartshorne, III.4, the Čech complex and
  Example~4.0.3, pp.~218--219.}
  Let \(C\), \(U = C \setminus \{P_0\}\), \(V = C \setminus \{Q_0\}\) be
  the two-affine cover of \cref{lem:exists_two_affine_cover}, and view
  \(\mathcal O_C\) in its \(\Module \bar k\)-carrier \(\toModuleKSheaf
  C\). For the two-element cover \(\mathfrak U = \{U, V\}\) the Čech
  complex of \(\mathcal O_C\) has only the two terms
  \[
    C^0 \;=\; \mathcal O_C(U) \oplus \mathcal O_C(V), \qquad
    C^1 \;=\; \mathcal O_C(U \cap V),
  \]
  with Čech differential
  \(d \colon C^0 \to C^1\),
  \(d(s, t) = s|_{U \cap V} - t|_{U \cap V}\) (the restriction of the
  section coming from \(U\) minus that coming from \(V\)). We define the
  \emph{first Čech cohomology of the two-affine cover} to be the
  cokernel of this single \(\bar k\)-linear map:
  \[
    \cechHI \;:=\;
    \operatorname{coker}\bigl(d \colon \mathcal O_C(U) \oplus
      \mathcal O_C(V) \longrightarrow \mathcal O_C(U \cap V)\bigr).
  \]
  Because the cover has only two members there is no higher term, so
  \(\cechHI\) is exactly the degree-\(1\) Čech cohomology
  \(\check H^1(\mathfrak U, \mathcal O_C)\).
\end{definition}

\begin{lemma}
[Čech equals derived cohomology on the two-affine cover]
  \label{lem:HModule_H1_iso_cech_two_affine}
  \lean{AlgebraicGeometry.Scheme.HModule_one_iso_cechH1}
  \uses{lem:exists_two_affine_cover, def:cech_H1_two_affine,
        def:Scheme_HModule}
  % SOURCE: [Hartshorne], III.4, Theorem~4.5 (Čech cohomology on an open
  % affine cover of a noetherian separated scheme agrees with derived
  % cohomology) (read from
  % references/hartshorne-algebraic-geometry.pdf, PDF page 239 = doc
  % p.222).
  % SOURCE QUOTE: "Theorem 4.5. \emph{Let $X$ be a noetherian separated
  % scheme, let $\mathfrak{U}$ be an open affine cover of $X$, and let
  % $\mathscr{F}$ be a quasi-coherent sheaf on $X$. Then for all $p \geq
  % 0$, the natural maps $(4.4)$ give isomorphisms}
  % $\check{H}^p(\mathfrak{U},\mathscr{F}) \xrightarrow{\sim}
  % H^p(X,\mathscr{F})$."
  \textit{Source: Hartshorne, III.4, Theorem~4.5, p.~222.}
  With the cover \(\mathfrak U = \{U, V\}\) of
  \cref{lem:exists_two_affine_cover}, the derived first cohomology of
  the structure sheaf is \(\bar k\)-linearly isomorphic to the Čech
  cokernel of \cref{def:cech_H1_two_affine}:
  \[
    H^1(C, \mathcal O_C) \;=\;
    \HModule\;\bar k\;(\toModuleKSheaf C)\;1 \;\cong\; \cechHI .
  \]
\end{lemma}

% SOURCE QUOTE PROOF: "PROOF. For $p = 0$ we have an isomorphism by
% (4.1). For the general case, embed $\mathscr{F}$ in a flasque,
% quasi-coherent sheaf $\mathscr{G}$ (3.6), and let $\mathscr{R}$ be the
% quotient: $0 \to \mathscr{F} \to \mathscr{G} \to \mathscr{R} \to 0$.
% For each $i_0 < \ldots < i_p$, the open set $U_{i_0,\ldots,i_p}$ is
% affine, since it is an intersection of affine open subsets of a
% separated scheme (II, Ex. 4.3). Since $\mathscr{F}$ is
% quasi-coherent, we therefore have an exact sequence ... Therefore we
% get its long exact sequence of Čech cohomology groups. Since
% $\mathscr{G}$ is flasque, its Čech cohomology vanishes for $p > 0$ by
% (4.3) ... we conclude that the natural map
% $\check{H}^1(\mathfrak{U},\mathscr{F}) \xrightarrow{\sim}
% H^1(X,\mathscr{F})$ is an isomorphism. But $\mathscr{R}$ is also
% quasi-coherent (II, 5.7), so we obtain the result for all $p$ by
% induction."
\begin{proof}
  \uses{lem:exists_two_affine_cover, def:cech_H1_two_affine,
        def:Scheme_HModule}
  A smooth proper curve over \(\bar k\) is projective (Hartshorne's
  curve convention, IV.1), hence a noetherian separated scheme, and the
  structure sheaf \(\mathcal O_C\) is quasi-coherent. By
  \cref{lem:exists_two_affine_cover} the cover \(\mathfrak U = \{U, V\}\)
  is an open affine cover of \(C\). Theorem~III.4.5 of Hartshorne then
  gives, for every \(p \ge 0\), a natural isomorphism
  \(\check H^p(\mathfrak U, \mathcal O_C) \xrightarrow{\;\sim\;}
  H^p(C, \mathcal O_C)\) between Čech and derived cohomology. Taking
  \(p = 1\) and identifying \(\check H^1(\mathfrak U, \mathcal O_C)\)
  with the explicit cokernel \(\cechHI\) of
  \cref{def:cech_H1_two_affine} (the two-element cover has no higher
  term) yields the stated \(\bar k\)-linear isomorphism
  \(H^1(C, \mathcal O_C) \cong \cechHI\).

  % NOTE: The Čech-to-derived comparison rests on the two-affine cover
  % being acyclic for the structure sheaf: higher cohomology of a quasi-coherent
  % sheaf over an affine open vanishes (the affine acyclicity input behind
  % Hartshorne III.4.5).
\end{proof}

\begin{lemma}
[The two-affine Čech cokernel is finite-dimensional]
  \label{lem:cech_H1_two_affine_finiteDimensional}
  \lean{AlgebraicGeometry.Scheme.cechH1_finiteDimensional}
  \uses{def:cech_H1_two_affine}
  The Čech cokernel \(\cechHI\) of \cref{def:cech_H1_two_affine} is a
  finite-dimensional \(\bar k\)-vector space:
  \(\FiniteDimensional\;\bar k\;\cechHI\), equivalently
  \(\Module.\Finite\;\bar k\;\cechHI\).
\end{lemma}

\begin{proof}
  \uses{def:cech_H1_two_affine}
  Write \(W = U \cap V = C \setminus \{P_0, Q_0\}\). The cokernel
  \(\cechHI = \mathcal O_C(W) / \bigl(\mathcal O_C(U) + \mathcal
  O_C(V)\bigr)\) is the obstruction to writing a regular function on
  \(W\) as a difference of a function regular on \(U\) and one regular
  on \(V\). A class in \(\cechHI\) is detected entirely by the finite
  \emph{polar parts} of its representative at the two excised points
  \(P_0\) and \(Q_0\): a representative whose principal parts at both
  points are realised by functions regular on the complementary affine
  is a coboundary. At each smooth closed point the space of admissible
  principal parts of order \(\le n\) is finite-dimensional, and the
  dimension count for a complete curve (Riemann's inequality) bounds the
  surviving polar-part classes by a quantity independent of \(n\). Hence
  only finitely many classes survive and \(\cechHI\) is
  finite-dimensional over \(\bar k\).

  % NOTE: Although framed as "the cokernel of one explicit \(\bar k\)-linear map",
  % this finite-dimensionality is genuine geometric content:
  % \(\dim_{\bar k} \cechHI\) is the genus of \(C\)
  % (cf. Hartshorne III.4, Example~4.0.3 and Exercise~4.7, which evaluate
  % exactly such cokernels for explicit curves to finite-dimensional spaces).
\end{proof}

\begin{lemma}
[Finiteness of \(H^1\) of the structure sheaf]
  \label{lem:HModule_structureSheaf_H1_finite}
  \lean{AlgebraicGeometry.Scheme.HModule_structureSheaf_H1_finite}
  \uses{lem:HModule_H1_iso_cech_two_affine,
        lem:cech_H1_two_affine_finiteDimensional}
  Let \(C\) be a smooth proper geometrically irreducible curve over
  \(\bar k\). Then \(H^1(C, \mathcal O_C) =
  \HModule\;\bar k\;(\toModuleKSheaf C)\;1\) is a finite-dimensional
  \(\bar k\)-vector space:
  \(\FiniteDimensional\;\bar k\;\bigl(\HModule\;\bar k\;
  (\toModuleKSheaf C)\;1\bigr)\). This is the real base obligation of
  the chapter, replacing any appeal to general Serre finiteness.
\end{lemma}

\begin{proof}
  \uses{lem:HModule_H1_iso_cech_two_affine,
        lem:cech_H1_two_affine_finiteDimensional}
  By \cref{lem:HModule_H1_iso_cech_two_affine} there is a
  \(\bar k\)-linear isomorphism \(H^1(C, \mathcal O_C) \cong \cechHI\).
  By \cref{lem:cech_H1_two_affine_finiteDimensional} the target
  \(\cechHI\) is a finite-dimensional \(\bar k\)-vector space.
  Finite-dimensionality transports across a linear isomorphism, so
  \(H^1(C, \mathcal O_C)\) is finite-dimensional over \(\bar k\).
\end{proof}

\section{The base case: finiteness for the structure sheaf}

\begin{lemma}
[Finiteness of the cohomology of the structure sheaf]
  \label{lem:HModule_structureSheaf_finite}
  \lean{AlgebraicGeometry.Scheme.HModule_structureSheaf_finite}
  \uses{lem:finrank_H0_toModuleKSheaf_eq_one,
        lem:HModule_structureSheaf_H1_finite,
        lem:grothendieck_vanishing_curve,
        def:Scheme_HModule}
  Let \(C\) be a smooth proper geometrically irreducible curve over
  \(\bar k\). Then the structure sheaf \(\mathcal O_C\) (in its
  \(\Module \bar k\)-flavoured carrier \(\toModuleKSheaf C\)) has finite
  cohomology: \(H^i(C, \mathcal O_C)\) is a finite-dimensional
  \(\bar k\)-vector space for every \(i\).

\end{lemma}

\begin{proof}
  \uses{lem:finrank_H0_toModuleKSheaf_eq_one,
        lem:HModule_structureSheaf_H1_finite,
        lem:grothendieck_vanishing_curve}
  By \cref{lem:grothendieck_vanishing_curve} it suffices to treat
  degrees \(0\) and \(1\).

  \emph{Degree \(0\).} On the proper geometrically integral curve
  \(C / \bar k\) the global functions are exactly the constants:
  \(\dim_{\bar k} H^0(C, \mathcal O_C) = 1\)
  (\cref{lem:finrank_H0_toModuleKSheaf_eq_one}). Hence
  \(H^0(C, \mathcal O_C) \cong \bar k\) is finite-dimensional.

  \emph{Degree \(1\).} The structure sheaf \(H^1\) is finite-dimensional
  by the two-affine Čech computation
  \cref{lem:HModule_structureSheaf_H1_finite}: \(H^1(C, \mathcal O_C)\)
  is isomorphic to the explicit Čech cokernel \(\cechHI\), which is a
  finite-dimensional \(\bar k\)-vector space. No appeal to general Serre
  finiteness is made; this base case is proven internally.

  Combining the two degrees with the degree-\(\ge 2\) vanishing,
  \(\mathcal O_C\) has finite cohomology.
\end{proof}

\section{Finiteness for every line bundle \(\mathcal O_C(D)\)}

\begin{theorem}
[Finiteness of the cohomology of \(\mathcal O_C(D)\) for every Weil divisor]
  \label{thm:HModule_lineBundle_finite}
  \lean{AlgebraicGeometry.Scheme.HModule_sheafOf_finite}
  \uses{lem:HModule_finite_of_ses, lem:HModule_skyscraper_finite,
        lem:HModule_structureSheaf_finite, lem:sheafOf_ses_single_add,
        lem:grothendieck_vanishing_curve, def:sheafOf,
        def:codim1_cycles}
  Let \(C\) be a smooth proper geometrically irreducible curve over
  \(\bar k\) and let \(D \in \mathrm{Div}(C)\) be any Weil divisor
  (\cref{def:codim1_cycles}). Then the invertible sheaf
  \(\mathcal O_C(D)\) of \cref{def:sheafOf} has finite cohomology:
  \(H^i(C, \mathcal O_C(D))\) is a finite-dimensional \(\bar k\)-vector
  space for every \(i\).

\end{theorem}

\begin{proof}
  \uses{lem:HModule_finite_of_ses, lem:HModule_skyscraper_finite,
        lem:HModule_structureSheaf_finite, lem:sheafOf_ses_single_add,
        lem:grothendieck_vanishing_curve}
  By \cref{lem:grothendieck_vanishing_curve} the groups
  \(H^i(C, \mathcal O_C(D))\) vanish for \(i \ge 2\), so finiteness
  only has to be established in degrees \(0\) and \(1\); we prove the
  full ``finite cohomology'' statement by induction on the divisor
  \(D\).

  \emph{Base case \(D = 0\).} The line bundle of the zero divisor is the
  structure sheaf, \(\mathcal O_C(0) = \mathcal O_C\)
  (\cref{lem:sheafOf_zero}). Its cohomology is finite by
  \cref{lem:HModule_structureSheaf_finite}.

  \emph{Inductive step: add or remove one closed point.} A Weil divisor
  is a finite \(\mathbb Z\)-combination of closed points
  (\cref{def:codim1_cycles}), so it is reached from \(0\) by finitely
  many moves \(D' \rightsquigarrow D' + [P]\) and their inverses, for
  closed points \(P \in C\). For each such move, the closed-point short
  exact sequence \cref{lem:sheafOf_ses_single_add} of
  \cref{chap:RiemannRoch_OcOfD} provides
  \[
    0 \;\longrightarrow\; \mathcal O_C(D')
    \;\longrightarrow\; \mathcal O_C(D' + [P])
    \;\longrightarrow\; k(P)
    \;\longrightarrow\; 0
  \]
  in the category of sheaves of \(\bar k\)-modules, whose cokernel is
  the closed-point skyscraper \(k(P)\). The skyscraper has finite
  cohomology (\cref{lem:HModule_skyscraper_finite}). Hence two of the
  three terms of this short exact sequence have finite cohomology:
  \begin{itemize}
    \item to pass from \(D'\) to \(D' + [P]\): the left term
      \(\mathcal O_C(D')\) (inductive hypothesis) and the right term
      \(k(P)\) are finite, so the middle term
      \(\mathcal O_C(D' + [P])\) is finite by
      \cref{lem:HModule_finite_of_ses};
    \item to pass from \(D' + [P]\) to \(D'\) (the inverse move): the
      middle term \(\mathcal O_C(D' + [P])\) (inductive hypothesis) and
      the right term \(k(P)\) are finite, so the left term
      \(\mathcal O_C(D')\) is finite by
      \cref{lem:HModule_finite_of_ses}.
  \end{itemize}
  Either direction of the 2-of-3 property therefore transports finite
  cohomology across a single-point change of the divisor. Starting from
  the finite base case \(D = 0\) and applying the step finitely many
  times reaches an arbitrary \(D \in \mathrm{Div}(C)\), so
  \(\mathcal O_C(D)\) has finite cohomology. Combined with the
  degree-\(\ge 2\) vanishing, \(H^i(C, \mathcal O_C(D))\) is a
  finite-dimensional \(\bar k\)-vector space in every degree.
\end{proof}

\section{Downstream finite-dimensionality instances}
\label{sec:SerreFiniteness_instances}

The finiteness theorem \cref{thm:HModule_lineBundle_finite} is the
producer of the \(\bar k\)-finite-dimensionality facts that the
Riemann--Roch chapters consume as typeclass instances. Each of the
following follows immediately by specialising
\cref{thm:HModule_lineBundle_finite} to the relevant sheaf and degree
(and, for the zero divisor, to \cref{lem:HModule_structureSheaf_finite});
no further mathematical content is involved.

\begin{itemize}
  \item \emph{The \(\ell\)-invariant is well defined.} The instance
    \(\Module.\mathtt{Finite}\;\bar k\;\bigl(H^0(C, \mathcal O_C(D))\bigr)\)
    --- equivalently
    \(\mathtt{FiniteDimensional}\;\bar k\;(\HModule\;\bar k\;
    (\mathtt{sheafOf}\;D)\;0)\) --- is the \(i = 0\) case of
    \cref{thm:HModule_lineBundle_finite}. It makes
    \(\ell(D) = \dim_{\bar k} H^0(C, \mathcal O_C(D))\)
    (\cref{def:l_invariant}) an honest natural number.
  \item \emph{The Euler characteristic is well defined.} The instances
    in degrees \(0\) and \(1\) make both terms of
    \(\chi(\mathcal O_C(D)) = \dim_{\bar k} H^0 - \dim_{\bar k} H^1\)
    (\cref{def:eulerChar_curve}) finite, so the difference is a
    well-defined integer.
  \item \emph{The genus is well defined.} The \(i = 1\) instance for the
    structure sheaf, \(\mathtt{FiniteDimensional}\;\bar k\;
    (\HModule\;\bar k\;(\toModuleKSheaf C)\;1)\), comes from
    \cref{lem:HModule_structureSheaf_finite}; it makes
    \(g(C) = \dim_{\bar k} H^1(C, \mathcal O_C)\) (\cref{def:genus}) a
    genuine natural number.
  \item \emph{The four \texttt{S2} side-conditions of the
    \(\chi\)-induction.} These four instances ---
    \(\bar k\)-finite-dimensionality of
    \(H^0(C, \mathcal O_C(D))\), \(H^1(C, \mathcal O_C(D))\),
    \(H^0(C, \mathcal O_C(D + [P]))\), and
    \(H^1(C, \mathcal O_C(D + [P]))\) --- are exactly the finiteness
    hypotheses appearing in the \(\chi\)-additive inductive step
    \cref{lem:euler_char_sheafOf_succ} of
    \cref{chap:RiemannRoch_RRFormula}, and are all discharged by
    \cref{thm:HModule_lineBundle_finite}.
\end{itemize}

\section{Out of scope}

The chapter is deliberately narrow. The following are \emph{not} part of
the \texttt{S2} build:
\begin{itemize}
  \item \emph{General Serre coherent-cohomology finiteness.} The full
    projective-dévissage / proper-pushforward proof of Serre's
    finiteness theorem is not formalised. Instead the chapter computes
    the single base case \(H^1(C, \mathcal O_C)\) directly from a
    two-affine Čech cover (\cref{lem:HModule_structureSheaf_H1_finite})
    and propagates finiteness from it by the formalised 2-of-3 and
    induction arguments. The chapter does not build a proper-pushforward
    or coherent-cohomology-finiteness theory.
  \item \emph{Computing the dimensions.} This chapter proves only
    \emph{finiteness}; the actual values \(\ell(D)\), \(\chi(\mathcal
    O_C(D)) = \deg D + 1 - g\), and the genus-\(0\) formula
    \(\ell(D) = \deg D + 1\) are the business of
    \cref{chap:RiemannRoch_RRFormula} and its
    \(\chi\)-identity, which consume the present finiteness instances.
\end{itemize}
